\chapter{Valuation}\label{ch:valuation}

% Scope: DCF, multiples (EV/EBITDA, P/E, EV/Revenue), terminal value, scenario weighting.
% Cross-link: the terminal-value formula is an application of the growing perpetuity --
%   \cref{eq:gordon} in \cref{ch:corporate-finance-core}.

\section{Discounted Cash Flow}\label{sec:dcf}
% complexity: 5 -> figure dcf-bridge

What is the entire firm worth today? DCF answers by converting every future free cash flow
--- across an explicit forecast horizon and a terminal value beyond it --- into a single
present value at the cost of capital. The answer determines whether to invest, acquire, or
sell. A firm is a bundle of future claims on cash: each year's free cash flow is discounted
individually, exactly as in \cref{eq:pv}; the terminal value captures all the cash flows
beyond the forecast horizon as a single lump. In practice, the terminal value almost always
dominates, which is why the growth assumptions embedded in it deserve more scrutiny than the
explicit-year forecasts.

Free cash flow is cash generated after taxes (levied on EBIT, which preserves the depreciation shield) and reinvestment needs are met:
\[
  \mathrm{FCF}_t \;=\; \mathrm{EBITDA}_t \;-\; \text{taxes}_t
                       \;-\; \Delta\mathrm{WC}_t \;-\; \mathrm{CapEx}_t.
\]
Discount each year's FCF at the WACC (\cref{eq:wacc}) and add the present value of the
terminal value \(\mathrm{TV}\) at the horizon year \(n\):
\begin{equation}\label{eq:dcf}
  EV \;=\; \sum_{t=1}^{n}\frac{\mathrm{FCF}_t}{(1+\mathrm{WACC})^{t}}
           \;+\; \frac{\mathrm{TV}}{(1+\mathrm{WACC})^{n}}.
\end{equation}
The WACC (\cref{eq:wacc}) is derived in \cref{ch:risk-and-capital-structure}; its equity
component is the CAPM cost of equity (\cref{eq:capm}). The terminal value formula is
developed in \cref{sec:terminal-value}. Note that \cref{eq:dcf} is the multi-period
extension of \cref{eq:npv} in \cref{ch:corporate-finance-core}: the initial outlay \(I_0\)
is the acquisition price, and enterprise value is the sum that makes \(\mathrm{NPV}=0\).

DCF is only as good as its inputs. WACC is estimated, not observed; FCF forecasts are
narratives dressed as numbers; and the terminal value --- typically
\qty{60}{\percent}--\qty{80}{\percent} of EV --- inherits every assumption in the explicit
period and then extrapolates them forever. Sensitivity to WACC is non-linear: a
one-percentage-point rise in the discount rate can cut enterprise value by
\qty{20}{\percent}--\qty{30}{\percent} in high-growth firms where most value lies in the
terminal period.

\begin{example}[Five-year DCF for the running B2B SaaS]
Five explicit forecast years for the running B2B SaaS. Year-1 FCF is \money{600000},
growing at \qty{30}{\percent} in years 1--3 (product-market-fit phase) then decelerating
to \qty{10}{\percent} in years 4--5. WACC is \qty{11}{\percent} (\cref{eq:wacc}
in \cref{ch:risk-and-capital-structure}; it equals the cost of equity here because
the business is equity-financed).

\[
\begin{array}{lrrrrrr}
  & \mathrm{Y1} & \mathrm{Y2} & \mathrm{Y3} & \mathrm{Y4} & \mathrm{Y5} \\[2pt]
  \mathrm{FCF} & \num{600000} & \num{780000} & \num{1014000} & \num{1115400} & \num{1226940} \\
  \text{Discount factor} & 0.901 & 0.812 & 0.731 & 0.659 & 0.593 \\
  \mathrm{PV} & \num{540541} & \num{633246} & \num{741379} & \num{735047} & \num{727612} \\
\end{array}
\]

Sum of explicit-period PVs: \(\approx\money{3377825}\).
Terminal value (year-5 FCF normalised at \qty{3}{\percent} steady-state growth;
computed fully in \cref{sec:terminal-value}): \(\approx\money{9400000}\) PV.
Enterprise value: \(\money{3377825}+\money{9400000}\approx\money{12800000}\).

Managerial implication: the explicit years contribute only \qty{26}{\percent} of
total EV. Raising the terminal growth assumption from \qty{3}{\percent} to
\qty{4}{\percent} adds over \money{1000000}; shaving WACC one point adds more.
Investors who fixate on near-term revenue miss where the value actually lives.
\end{example}

\begin{figure}[htbp]
  \centering
  \includegraphics[width=0.86\linewidth]{dcf-bridge}
  \caption{DCF waterfall: PV of each explicit-year FCF stacked against PV of terminal
  value. Terminal value dominates (\(\approx\)\qty{74}{\percent} of enterprise value),
  illustrating why growth-rate and discount-rate assumptions in the terminal period
  outweigh near-term forecast accuracy.}
  \label{fig:dcf-bridge}
\end{figure}

\begin{admonition}{Warning}
Treating the DCF output as \emph{the} value rather than one scenario. The model
answers ``what is value \emph{given these assumptions}?'' --- not ``what will a
buyer pay?'' A precise DCF built on optimistic growth and a low WACC is a
sophisticated way to confirm a prior. The discipline is running the model with
assumptions you would defend in a credit committee, not assumptions you need to
make the deal work.
\end{admonition}

Levered DCF discounts equity free cash flows at the cost of equity directly; unlevered DCF
discounts firm FCFs at the unlevered cost and then deducts debt. Adjusted present value
(APV) separates the unlevered firm value from the PV of the tax shield, useful when
leverage changes materially over the forecast horizon --- an LBO or a staged infrastructure
build, for instance.

DCF formalises the same logic as the NPV rule in \cref{ch:corporate-finance-core}
(\cref{eq:npv}) and the CLV formula in \cref{ch:customer-economics} (\cref{eq:clv}):
all are present values of future cash streams, discounted at the appropriate cost
of capital. \textcite{brealey2020}, \textcite{damodaran2012}, and
\textcite{berk6thcorporate} each develop the DCF framework; \textcite{damodaran2012}
is the most rigorous treatment of terminal-value mechanics.

\section{Terminal Value}\label{sec:terminal-value}
% complexity: 4

Beyond the explicit forecast horizon, the firm does not cease to exist. Terminal value
converts post-horizon cash flows into a single number anchored at year \(n\), without
extending a twenty-year spreadsheet. At the horizon the firm is assumed to reach a
steady state: growth stabilises at a long-run rate \(g\), margins stop expanding. The
growing-perpetuity formula (\cref{eq:gordon}) prices that steady state; discounting it
back to today gives the terminal value's present contribution to enterprise value.

Let \(\mathrm{FCF}_{n+1}\) be the first post-horizon cash flow (year \(n\) FCF grown by
one period) and \(g\) the perpetuity growth rate. The terminal value at year \(n\) is a
direct application of \cref{eq:gordon}:
\begin{equation}\label{eq:terminal-value}
  \mathrm{TV} \;=\; \frac{\mathrm{FCF}_{n+1}}{\mathrm{WACC}-g}.
\end{equation}
This is why \cref{eq:terminal-value} is already implicit in every DCF: plug
\(\mathrm{TV}/(1+\mathrm{WACC})^n\) into \cref{eq:dcf} and the full formula collapses to
a two-stage growing perpetuity. Brand strength raises \(g\) at the horizon
(\cref{ch:brand-growth}), compounding through \cref{eq:terminal-value} into enterprise
value.

Terminal value is mechanically correct but economically fragile. \(g\) must be below WACC
(the convergence condition of \cref{eq:gordon}); it must also be below long-run nominal
GDP growth, since no firm can outgrow the whole economy forever. Pinning
\(g=\qty{5}{\percent}\) when WACC is \qty{11}{\percent} is defensible; pinning
\(g=\qty{9}{\percent}\) makes the denominator near-zero and TV blows up.

\begin{example}[Terminal value for the running SaaS]
Year-5 normalised FCF for the SaaS is \money{1226940}; applying one period of
steady-state growth (\qty{3}{\percent}) gives \(\mathrm{FCF}_{6}=\money{1263748}\).
With WACC at \qty{11}{\percent}:
\[
  \mathrm{TV} \;=\; \frac{\money{1263748}}{0.11-0.03} \;=\; \money{15796850}.
\]
Discounting that year-5 terminal value back five years at \qty{11}{\percent} gives its
present contribution: \(\mathrm{PV}(\mathrm{TV})=\money{15796850}/1.11^5\approx\money{9400000}\).
Added to the \money{3377825} of explicit-period PVs, enterprise value is
\(\approx\money{12800000}\); the terminal value is \qty{74}{\percent} of it ---
confirming that the valuation is largely a bet on the steady state, not the next five
years.
\end{example}

\begin{admonition}{Warning}
\emph{Double-counting growth} is the canonical error: building high growth into
the explicit period \emph{and} the terminal perpetuity. The terminal \(g\) must
reflect only the period \emph{after} the explicit years, once competitive
dynamics have compressed margins toward the industry steady state. Reusing the
same \qty{30}{\percent} growth rate from year 1 in the terminal value produces an
absurdity. The same trap is flagged for \cref{eq:gordon} itself in
\cref{sec:growing-perpetuity}.
\end{admonition}

An alternative terminal value methodology is the exit-multiple approach: estimate the
EBITDA multiple a buyer would pay at the horizon and multiply by year-\(n\) EBITDA.
Importantly, this is not independent of DCF --- a multiple is itself a compressed
growing-perpetuity (\cref{sec:multiples}), so the two methods converge when their embedded
growth and discount assumptions are consistent. Using both as a cross-check, not as
independent valuations, is best practice.

Terminal value is the dominant component of enterprise value for any firm whose principal
value lies beyond the explicit forecast horizon --- that is, most growth-stage companies.
\textcite{damodaran2012} provides the most systematic treatment; \textcite{brealey2020}
covers the consistency conditions between the exit-multiple and perpetuity approaches.

\section{Multiples}\label{sec:multiples}
% complexity: 3

You want a quick market-anchored value without a full DCF build. Multiples offer a shortcut
--- but only if you understand what growth and return assumptions they are compressing.
Misread the compression and you pay the wrong price. An EV/Revenue multiple looks like a
market opinion, but it is a DCF in disguise: the multiple embeds assumptions about margins,
growth, and discount rates. Understanding the algebra lets you audit whether a traded
multiple is cheap or expensive relative to fundamentals, and why multiples collapsed from
2021 to 2023.

Start from \cref{eq:gordon}: \(EV = \mathrm{FCF}/(r-g)\). Express FCF as a margin on
revenue \(R\): \(\mathrm{FCF}=m\cdot R\). Divide both sides by \(R\) and allow one period
of growth so \(\mathrm{FCF}_1 = m\cdot R\cdot(1+g)\):
\begin{equation}\label{eq:ev-revenue}
  \frac{EV}{R} \;=\; \frac{m\,(1+g)}{\mathrm{WACC}-g}.
\end{equation}
This is the DCF-consistent revenue multiple. It rises with margin \(m\), rises with growth
\(g\), and \emph{falls with WACC} --- which is why rising interest rates compress multiples
mechanically, independent of any change in the underlying business. \cref{eq:ev-revenue}
assumes steady-state margins and growth; traded multiples also embed sentiment, liquidity
premia, and peer-group herding. A multiple without a view on \(m\), \(g\), and WACC is not
a valuation --- it is a price, which may or may not reflect value.

\begin{example}[Revenue multiple versus DCF for the running SaaS]
The SaaS has \qty{20}{\percent} FCF margin, \(g=\qty{3}{\percent}\), WACC \qty{11}{\percent}.
From \cref{eq:ev-revenue}:
\[
  \frac{EV}{R} \;=\; \frac{0.20\times1.03}{0.11-0.03} \;=\; \frac{0.206}{0.08} \;=\; 2.6\times.
\]
With ARR of \money{4000000}:
\[
  EV \;=\; 2.6\times\money{4000000} \;=\; \money{10400000}.
\]
The DCF in \cref{sec:dcf} yields \(\approx\money{12800000}\) --- the
\money{2400000} gap is the \emph{growth premium} that steady-state multiples miss:
the model uses only \(g=\qty{3}{\percent}\) (long-run), but the firm is currently
growing at \qty{30}{\percent}. A multiple calibrated to peers growing at
\qty{30}{\percent} would clear far higher. The discipline is understanding which
growth rate is embedded before citing a multiple as evidence of value.
\end{example}

In 2021, public SaaS traded at 20--30$\times$ revenue; by late 2023 multiples had
compressed to 5--8$\times$. The dominant driver was not deteriorating unit economics but a
mechanical re-pricing: the risk-free rate rose from \qty{0}{\percent} to \qty{5}{\percent},
lifting WACC and collapsing the denominator in \cref{eq:ev-revenue}. Whether 2021 multiples
reflected irrational exuberance or 2023 multiples over-corrected is contested; what is not
contested is that \cref{eq:ev-revenue} tracks the direction and approximate magnitude of the
compression (\cref{ch:risk-and-capital-structure}).

The EV/EBITDA multiple strips out capital structure and depreciation policy, making
cross-company comparison cleaner than P/E. Its DCF analogue replaces \(m\cdot R\) with
EBITDA\(\times(1-t)\times(1-\mathrm{reinvestment\,rate})\) --- economically identical once
tax and reinvestment are consistent. The practical advantage of multiples is speed and
market-anchoring; the risk is confusing a market price with a model of intrinsic value.

Multiples are best treated as a sanity-check on a DCF, not a replacement for one.
\textcite{damodaran2012} systematically decompose the algebra behind every common
multiple; the parallel with the Gordon model (\cref{eq:gordon}) is made explicit
in his discussion of the PEG ratio and the dividend-yield inversion.

\section{Scenario Analysis}\label{sec:scenario-analysis}
% complexity: 3

A point-estimate DCF or multiple implies certainty about an uncertain future. Scenario
analysis replaces that illusion with a probability-weighted distribution of enterprise
values, making the risk explicit and the decision defensible. Rather than asking ``what is
the value?'' scenario analysis asks ``what is the value \emph{conditional on each plausible
future}?'' Weight the outcomes by their probabilities and the expected value captures both
the upside and the downside in a single actionable number.

Let \(p_i\) be the probability of scenario \(i\) and \(EV_i\) the enterprise value under
that scenario (computed via \cref{eq:dcf} or a multiple from \cref{eq:ev-revenue}). The
probability-weighted enterprise value is:
\begin{equation}\label{eq:scenario-ev}
  EV^* \;=\; \sum_{i} p_i \cdot EV_i.
\end{equation}
\cref{eq:scenario-ev} reuses the expected-value machinery of \cref{eq:ev} in
\cref{ch:decisions-under-uncertainty}: it is the same probability-weighting applied to
valuations rather than payoffs. Under that chapter's framework, a risk-averse
decision-maker would further adjust downward for variance; here the adjustment is embedded
in the discount rate rather than the probabilities. The output depends entirely on the
quality of the \(p_i\) and \(EV_i\) estimates. Scenario probabilities are rarely objective
--- they reflect the analyst's priors. The formula disciplines the \emph{structure} of
thinking (three futures, assign weights, compute) but cannot substitute for sound judgement
about which futures are plausible.

\begin{example}[Scenario-weighted acquisition decision]
Three scenarios for the SaaS at a prospective acquisition decision:

\begin{description}
  \item[Bear (\(p=0.25\))] Competition intensifies; growth stalls; EV \money{8000000}.
  \item[Base (\(p=0.50\))] Current trajectory holds; EV \money{15000000}.
  \item[Bull (\(p=0.25\))] Platform network effects kick in; EV \money{28000000}.
\end{description}

Applying \cref{eq:scenario-ev}:
\[
  \begin{aligned}
    EV^* &= 0.25\times\money{8000000}
          + 0.50\times\money{15000000}
          + 0.25\times\money{28000000} \\
         &= \money{2000000}+\money{7500000}+\money{7000000}
          = \money{16500000}.
  \end{aligned}
\]
The expected EV is \money{16500000} --- above the single-point DCF of
\money{12800000} because the bull scenario is weighted equally with the bear yet sits
farther from the base.
The asymmetry between scenarios drives the result; sensitivity to the \(p_i\)
assignments matters more than the model precision within each scenario.
\end{example}

\begin{admonition}{Warning}
Collapsing three scenarios into a single ``base case'' to avoid the discomfort of
the downside. \emph{Probability weighting is not pessimism} --- it is accounting
for the full range of outcomes. A bear scenario with \(p=0.25\) is not a worst
case to be dismissed; it is a claim that one in four similar decisions ends there.
\end{admonition}

Monte Carlo simulation extends scenario analysis by drawing from continuous probability
distributions over each input (WACC, \(g\), year-1 FCF) and computing the full EV
distribution. The output is a histogram of enterprise values rather than three point
estimates, revealing fat tails that scenario weights can obscure. The structural question
--- which inputs drive variance --- is best answered by a Sobol sensitivity analysis, which
partitions total EV variance to each input.

Scenario analysis links back to the decision framework in \cref{ch:decisions-under-uncertainty}
and forward to the capital-structure choice in \cref{ch:risk-and-capital-structure}:
higher EV variance argues for lower leverage, since distress costs are more likely.
\textcite{damodaran2012} develops the simulation extension; \textcite{brealey2020}
covers the scenario-break-even approach as a complement.
